% !TeX root = ../main.tex
\chapter{Conclusions and Future Work}\label{ch:conclusions}

\begin{guidance}
Summarise the work done and the main results, and state clearly whether the objectives
set in \cref{ch:objectives} (those of the thesis proposal approved by the academic
committee) have been achieved. Do not introduce new results here. Finish with the
lines of future work that derive from the thesis. Some universities also ask for a
personal assessment of the competences acquired.
\end{guidance}

\section{Conclusions}\label{sec:conclusions}

[Main conclusions of the thesis.]

\section{Fulfilment of the objectives}\label{sec:fulfilment}

\Cref{tab:fulfilment} relates each specific objective to the part of the thesis where it
has been addressed.

\begin{table}[htbp]
  \centering
  \caption{Fulfilment of the specific objectives of the thesis.}
  \label{tab:fulfilment}
  \begin{tabularx}{\textwidth}{@{}l X l@{}}
    \toprule
    \textbf{Objective} & \textbf{How it has been achieved} & \textbf{Where} \\
    \midrule
    \ref{obj:1} & [Brief description] & \cref{ch:background} \\
    \ref{obj:2} & [Brief description] & \cref{sec:data} \\
    \ref{obj:3} & [Brief description] & \cref{sec:methods} \\
    \ref{obj:4} & [Brief description] & \cref{ch:results} \\
    \bottomrule
  \end{tabularx}
\end{table}

\section{Future work}\label{sec:future-work}

[Possible improvements and new lines of research that derive from this work.]
